\documentclass[journal]{IEEEtran}

\usepackage[utf8]{inputenc}
\usepackage[T1]{fontenc}
\usepackage{graphicx}
\usepackage{amsmath}
\usepackage{array}

\graphicspath{{./}{./fig/}}

\hyphenation{op-tical net-works semi-conduc-tor}


\begin{document}

\title{Experimento 08: Respuesta en Frecuencia de un Amplificador}
\author{Brandon~Najera~Zuñiga,~\IEEEmembership{Estudiante,~ITCR}
        y~Kevin~Mora~Sobalvarr.~\IEEEmembership{Estudiante,~ITCR,}
}
\markboth{EL3215 Laboratorio de Electr\'onica Anal\'ogica, IS2026}%
{EL3215 Laboratorio de Electr\'onica Anal\'ogica}

\maketitle

% ---------------------------------------------------------------------------
\begin{abstract}
Se estudi\'o la respuesta en frecuencia de un amplificador BJT en configuraci\'on de emisor com\'un construido con un transistor 2N3904 y alimentado con $V_{CC} = +15\,\text{V}$. En la Parte~1 se determinaron las tres frecuencias cr\'iticas de corte inferior asociadas a los capacitores de acople $C_1$, $C_2$ y $C_3$ mediante c\'alculo te\'orico y medici\'on experimental, siendo $C_2$ el capacitor dominante con una frecuencia cr\'itica calculada de $70{,}82\,\text{Hz}$ y medida de $80{,}87\,\text{Hz}$. Adem\'as, se calcul\'o el valor de $C_2$ necesario para fijar la frecuencia cr\'itica inferior general en $300\,\text{Hz}$, obteniendo $C_2 = 33{,}11\,\mu\text{F}$. En la Parte~2 se analiz\'o la respuesta en alta frecuencia mediante el modelo de Miller con capacitores externos de $100\,\text{pF}$, obteni\'endose una frecuencia cr\'itica superior general calculada de $140{,}55\,\text{kHz}$ y medida de $139{,}8\,\text{kHz}$, con un error del $0{,}54\,\%$, validando el modelo te\'orico de primer orden.
\end{abstract}

\begin{IEEEkeywords}
Respuesta en frecuencia, amplificador BJT, emisor com\'un, frecuencia cr\'itica, efecto Miller, capacitor de acople, 2N3904.
\end{IEEEkeywords}


% ===========================================================================
\section{Introducci\'on}

\IEEEPARstart{E}{n} un amplificador con transistores t\'ipico, la respuesta en baja frecuencia es determinada por los capacitores de acople del circuito. Para su an\'alisis se traza la ruta de descarga de cada capacitor, formando un circuito equivalente RC simplificado \cite{Floyd2008}. La respuesta en alta frecuencia est\'a dominada por el efecto Miller capacitivo asociado a la capacitancia interna base-colector del transistor.

En la Parte~1 se analizaron las tres frecuencias cr\'iticas inferiores del amplificador de la Figura~\ref{fig:circuito_p1} y se calcul\'o el valor de $C_2$ necesario para fijar la frecuencia de corte inferior en $300\,\text{Hz}$. En la Parte~2 se extendi\'o el an\'alisis al dominio de alta frecuencia mediante el modelo de Miller, utilizando capacitores externos $C_4 = C_5 = C_6 = 100\,\text{pF}$ para reducir las frecuencias cr\'iticas superiores a rangos medibles.

Los resultados obtenidos permiten verificar la validez de los modelos te\'oricos de primer orden para la estimaci\'on de frecuencias cr\'iticas en amplificadores BJT de una etapa. Sin embargo, en la pr\'actica se observaron algunas desviaciones que se discutir\'an posteriormente.


% ===========================================================================
\section{Parte 1: Respuesta en Baja Frecuencia}

\subsection{Circuito implementado}

El circuito empleado corresponde al amplificador BJT en emisor com\'un de la Figura~\ref{fig:circuito_p1}. Los resistores $R_A$ y $R_B$ conforman una etapa atenuadora de entrada para generar $V_{ent} = 20\,\text{mV}_{pp}$.

\begin{figure}[!ht]
\centering
\includegraphics[width=3.2in]{F1}
\caption{Amplificador en emisor com\'un --- Parte~1.}
\label{fig:circuito_p1}
\end{figure}

\begin{figure}[!ht]
\centering
\includegraphics[width=3.2in]{F1sim}
\caption{Circuito de la Parte~1 simulado en NI Multisim.}
\label{fig:circuito_p1sim}
\end{figure}


% ---------------------------------------------------------------------------
\subsection{Medici\'on de resistencias}

La Tabla~\ref{tab:resistencias} presenta los valores nominales y medidos de las resistencias utilizadas. Todos los valores se encuentran dentro de la tolerancia est\'andar del $5\,\%$.

\begin{table}[!ht]
    \centering
    \caption{Valores de resistencias utilizadas}
    \begin{tabular}{|>{\centering\arraybackslash}m{1.3cm}|
                     >{\centering\arraybackslash}m{2.5cm}|
                     >{\centering\arraybackslash}m{2.5cm}|}
        \hline
        Resistor & Valor nominal & Valor medido \\
        \hline
        $R_A$    & $1{,}0\,\text{k}\Omega$  & $0{,}989\,\text{k}\Omega$ \\
        \hline
        $R_B$    & $47\,\Omega$              & $46{,}75\,\Omega$          \\
        \hline
        $R_1$    & $68\,\text{k}\Omega$      & $67{,}68\,\text{k}\Omega$  \\
        \hline
        $R_2$    & $10\,\text{k}\Omega$      & $9{,}849\,\text{k}\Omega$  \\
        \hline
        $R_{E1}$ & $10\,\Omega$              & $9{,}912\,\Omega$          \\
        \hline
        $R_{E2}$ & $560\,\Omega$             & $562{,}22\,\Omega$         \\
        \hline
        $R_C$    & $3{,}9\,\text{k}\Omega$   & $3{,}819\,\text{k}\Omega$  \\
        \hline
        $R_L$    & $10\,\text{k}\Omega$      & $9{,}921\,\text{k}\Omega$  \\
        \hline
    \end{tabular}
    \label{tab:resistencias}
\end{table}

Los valores medidos de todos los resistores presentan desviaciones dentro del rango de tolerancia est\'andar del $5\,\%$. Estas peque\~nas diferencias afectan proporcionalmente los puntos de operaci\'on en CD y los par\'ametros de peque\~na se\~nal calculados.


% ---------------------------------------------------------------------------
\subsection{Par\'ametros CD y CA}

Los par\'ametros de polarizaci\'on en CD y de peque\~na se\~nal en CA se calcularon a partir de las resistencias medidas asumiendo $\beta_{ca} = 200$ y $V_{BE} = 0{,}7\,\text{V}$ \cite{Floyd2008}. Las expresiones empleadas se presentan en el Ap\'endice~\ref{apx:ec_p1}.

\begin{table}[!ht]
    \centering
    \caption{Par\'ametros CD y CA del amplificador --- Parte~1}
    \begin{tabular}{|>{\centering\arraybackslash}m{1.5cm}|
                     >{\centering\arraybackslash}m{1.5cm}|
                     >{\centering\arraybackslash}m{1.5cm}|
                     >{\centering\arraybackslash}m{1.5cm}|}
        \hline
        Par\'ametro & Calculado & Simulado & Medido \\
        \hline
        $V_B$ (V)            & $1{,}906$ & $1{,}923$ & $1{,}794$ \\
        \hline
        $V_E$ (V)            & $1{,}206$ & $1{,}123$ & $1{,}108$ \\
        \hline
        $I_E$ (mA)           & $2{,}107$ & $1{,}972$ & --- \\
        \hline
        $V_C$ (V)            & $6{,}953$ & $7{,}316$ & $7{,}636$ \\
        \hline
        $V_{CE}$ (V)         & $5{,}747$ & $6{,}193$ & $6{,}535$ \\
        \hline
        $r'_e$ $(\Omega)$    & $11{,}86$ & $13{,}19$ & --- \\
        \hline
        $A_v$                & $-126{,}63$ & $-120{,}98$ & $-215$ \\
        \hline
        $V_{out}$ (V$_{pp}$) & $2{,}533$ & $2{,}420$ & $4{,}400$ \\
        \hline
    \end{tabular}
    \label{tab:cd_ca_p1}
\end{table}

Los valores de tensi\'on de polarizaci\'on presentan errores entre $6\,\%$ y $12\,\%$ respecto a los medidos. El transistor opera en la regi\'on activa, confirmado por $V_{CE} = 6{,}535\,\text{V}$, comprendido entre $0{,}2\,\text{V}$ y $V_{CC}$. La discrepancia m\'as significativa se observa en $A_v$ y $V_{out}$, atribuible principalmente a que la suposici\'on $\beta_{ca} = 200$ puede diferir del $\beta$ real del transistor, y a efectos de carga del osciloscopio no contemplados en el modelo simplificado \cite{Floyd2008}. La se\~nal de salida simulada se muestra en la Figura~\ref{fig:vout1}.

\begin{figure}[!ht]
\centering
\includegraphics[width=3.2in]{VOUT1}
\caption{Se\~nal de salida $V_{out}$ y entrada $V_s$ simuladas --- Parte~1 (NI Multisim).}
\label{fig:vout1}
\end{figure}


% ---------------------------------------------------------------------------
\subsection{Resistencias equivalentes}

Las resistencias equivalentes se calcularon trazando la ruta de carga y descarga en CA para cada capacitor, con los dem\'as considerados cortocircuitos \cite{Floyd2008}:

\begin{equation}
    R_{eq}(C_1) = \beta_{ca}(R_{E1} + r'_e)\|R_1\|R_2 + (R_A + R_{th})\|R_B
    \label{eq:req_c1}
\end{equation}

\begin{equation}
    R_{eq}(C_2) = R_{ent(emisor)}\|R_{E2}
    \label{eq:req_c2}
\end{equation}

donde $R_{ent(emisor)} = r'_e + R_{E1} + R_{umbral}/\beta_{ca}$ con $R_{umbral} = R_A\|R_B\|R_1\|R_2$.

\begin{equation}
    R_{eq}(C_3) = R_C + R_L
    \label{eq:req_c3}
\end{equation}

\begin{table}[!ht]
    \centering
    \caption{Resistencias equivalentes para cada capacitor de acople}
    \begin{tabular}{|>{\centering\arraybackslash}m{2cm}|
                     >{\centering\arraybackslash}m{4cm}|}
        \hline
        Capacitor & $R_{eq}$ calculada \\
        \hline
        $C_1$ & $4{,}67\,\text{k}\Omega$ \\
        \hline
        $C_2$ & $22{,}5\,\Omega$          \\
        \hline
        $C_3$ & $13{,}9\,\text{k}\Omega$  \\
        \hline
    \end{tabular}
    \label{tab:req}
\end{table}

La resistencia equivalente de $C_2$ ($22{,}5\,\Omega$) es significativamente menor que las de $C_1$ y $C_3$, lo que anticipa que $C_2$ dominar\'a la respuesta en baja frecuencia, pues $f \propto 1/(R_{eq} C)$ \cite{Floyd2008}.


% ---------------------------------------------------------------------------
\subsection{Frecuencias cr\'iticas de corte inferior}

Las frecuencias cr\'iticas inferiores se calcularon con:

\begin{equation}
    f = \frac{1}{2\pi R_{eq} C}
    \label{eq:fc}
\end{equation}

La frecuencia cr\'itica inferior general se estim\'o como la suma de las frecuencias individuales, expresi\'on que provee una cota superior del valor real \cite{Floyd2008}:

\begin{equation}
    f_{cl} \approx f_{C_1} + f_{C_2} + f_{C_3}
    \label{eq:fcl}
\end{equation}

Las frecuencias individuales se midieron aislando cada capacitor con capacitores de $1000\,\mu\text{F}$ en paralelo con los otros dos, reduciendo la frecuencia hasta que la salida cayera al $70{,}7\,\%$ de su valor en banda media.

\begin{table}[!ht]
    \centering
    \caption{Frecuencias cr\'iticas de corte inferior}
    \begin{tabular}{|>{\centering\arraybackslash}m{1.4cm}|
                     >{\centering\arraybackslash}m{1.7cm}|
                     >{\centering\arraybackslash}m{1.7cm}|
                     >{\centering\arraybackslash}m{1.5cm}|}
        \hline
        Capacitor & Calculada (Hz) & Medida (Hz) & Simulada (Hz) \\
        \hline
        $C_1$            & $34{,}06$ & $50{,}98$ & $50{,}9$  \\
        \hline
        $C_2$            & $70{,}82$ & $80{,}87$ & $61{,}5$  \\
        \hline
        $C_3$            & $52{,}04$ & $30{,}01$ & $66{,}5$  \\
        \hline
        General $f_{cl}$ & $156{,}92$ & --- & --- \\
        \hline
    \end{tabular}
    \label{tab:fc_inf}
\end{table}

Los resultados muestran discrepancias entre los tres m\'etodos. Para $C_1$, el valor medido y simulado coinciden bien entre s\'i pero superan el calculado, lo que sugiere que la resistencia equivalente real vista por $C_1$ es menor a la te\'orica, posiblemente por una resistencia interna de la fuente no contemplada en el modelo. Para $C_2$, los tres valores son los m\'as cercanos entre s\'i. La mayor discrepancia se presenta en $C_3$: el valor calculado y simulado difieren significativamente del medido, atribuible a resistencias par\'asitas del protoboard que incrementan la resistencia equivalente efectiva y desplazan la frecuencia cr\'itica hacia valores menores. El capacitor dominante es $C_2$, con la mayor frecuencia cr\'itica en medici\'on ($80{,}87\,\text{Hz}$). Las mediciones individuales se ilustran en las Figuras~\ref{fig:fre1}, \ref{fig:fre2} y~\ref{fig:fre3}.

\begin{figure}[!ht]
\centering
\includegraphics[width=3.2in]{Fre1}
\caption{Medici\'on de la frecuencia cr\'itica de $C_1$ ($f_{C_1} = 50{,}98\,\text{Hz}$).}
\label{fig:fre1}
\end{figure}

\begin{figure}[!ht]
\centering
\includegraphics[width=3.2in]{Fre2}
\caption{Medici\'on de la frecuencia cr\'itica de $C_2$ ($f_{C_2} = 80{,}87\,\text{Hz}$).}
\label{fig:fre2}
\end{figure}

\begin{figure}[!ht]
\centering
\includegraphics[width=3.2in]{Fre3}
\caption{Medici\'on de la frecuencia cr\'itica de $C_3$ ($f_{C_3} = 30{,}01\,\text{Hz}$).}
\label{fig:fre3}
\end{figure}


% ---------------------------------------------------------------------------
\subsection{Ajuste de frecuencia cr\'itica a 300 Hz}

Para fijar $f_{cl} = 300\,\text{Hz}$ se asign\'o la contribuci\'on restante a $C_2$, descontadas las aportaciones de $C_1$ y $C_3$:

\begin{equation}
    f_{C_2,\text{nuevo}} = 300 - f_{C_1} - f_{C_3}
\end{equation}

\begin{equation}
    C_2 = \frac{1}{2\pi \cdot R_{eq}(C_2) \cdot f_{C_2,\text{nuevo}}}
\end{equation}

Sustituyendo con los valores calculados de $f_{C_1} = 34{,}06\,\text{Hz}$, $f_{C_3} = 52{,}04\,\text{Hz}$ y $R_{eq}(C_2) = 22{,}47\,\Omega$:

\begin{equation}
    f_{C_2,\text{nuevo}} = 300 - 34{,}06 - 52{,}04 = 213{,}90\,\text{Hz}
\end{equation}

\begin{equation}
    C_2 = \frac{1}{2\pi \cdot 22{,}47\,\Omega \cdot 213{,}90\,\text{Hz}} = 33{,}11\,\mu\text{F}
\end{equation}

\begin{table}[!ht]
    \centering
    \caption{Ajuste de la frecuencia cr\'itica inferior mediante $C_2$}
    \begin{tabular}{|>{\centering\arraybackslash}m{3.5cm}|
                     >{\centering\arraybackslash}m{3.5cm}|}
        \hline
        $C_2$ calculado & $f_{cl}$ medida con nuevo $C_2$ \\
        \hline
        $33{,}11\,\mu\text{F}$ & $200\,\text{Hz}$ \\
        \hline
    \end{tabular}
    \label{tab:ajuste_fc}
\end{table}

El valor comercial m\'as cercano es $33\,\mu\text{F}$.

% ---------------------------------------------------------------------------
\subsection{Preguntas de la Parte 1}

\begin{enumerate}

\item \textbf{¿Cu\'al capacitor tiene mayor impacto en la frecuencia cr\'itica inferior? Justifique.}

El capacitor $C_2$ presenta el mayor impacto. Su frecuencia cr\'itica medida es la mayor de las tres ($80{,}87\,\text{Hz}$), frente a $50{,}98\,\text{Hz}$ de $C_1$ y $30{,}01\,\text{Hz}$ de $C_3$. Esto se justifica mediante la ecuaci\'on~(\ref{eq:fc}): la resistencia equivalente vista por $C_2$ es de solo $22{,}47\,\Omega$, significativamente menor que la de $C_1$ ($4{,}672\,\text{k}\Omega$) y $C_3$ ($13{,}90\,\text{k}\Omega$). Dado que $f \propto 1/(R_{eq} C)$, la menor resistencia de $C_2$ eleva su frecuencia cr\'itica a pesar de contar con la mayor capacitancia ($100\,\mu\text{F}$) \cite{Floyd2008}.

\item \textbf{¿Qu\'e cambio se podr\'ia hacer para reducir la frecuencia de corte inferior por un factor de 5?}

De la ecuaci\'on~(\ref{eq:fc}), la frecuencia es inversamente proporcional al producto $R_{eq} \cdot C$. Para reducir $f_{cl}$ por un factor de 5 sin modificar la topolog\'ia del circuito, la alternativa m\'as pr\'actica consiste en incrementar los capacitores de acople en ese mismo factor. En particular, el capacitor dominante $C_2$ deber\'ia aumentarse de $100\,\mu\text{F}$ a $500\,\mu\text{F}$, o bien multiplicar los tres capacitores por 5 simult\'aneamente \cite{Floyd2008}.

\end{enumerate}


% ===========================================================================
\section{Parte 2: Respuesta en Alta Frecuencia}

\subsection{Circuito implementado}

Se utiliz\'o el mismo amplificador de la Parte~1 a\~nadiendo $C_4 = C_5 = C_6 = 100\,\text{pF}$ entre las terminales del transistor seg\'un la Figura~\ref{fig:circuito_p2}, reduciendo las frecuencias cr\'iticas superiores a rangos medibles sin alterar el punto de operaci\'on en CD.

\begin{figure}[!ht]
\centering
\includegraphics[width=3.2in]{F2}
\caption{Amplificador en emisor com\'un para pruebas en alta frecuencia ($C_4=C_5=C_6=100\,\text{pF}$).}
\label{fig:circuito_p2}
\end{figure}

\begin{figure}[!ht]
\centering
\includegraphics[width=3.2in]{F2sim}
\caption{Circuito de la Parte~2 simulado en NI Multisim.}
\label{fig:circuito_p2sim}
\end{figure}


% ---------------------------------------------------------------------------
\subsection{Par\'ametros CD y CA}

Los capacitores $C_4$, $C_5$ y $C_6$ son circuitos abiertos en CD, por lo que el punto de operaci\'on es id\'entico al de la Parte~1. Los valores de $V_E$, $V_C$ y $V_{CE}$ en CD se verificaron en simulaci\'on (Figuras~\ref{fig:ve2}, \ref{fig:vc2} y~\ref{fig:vce2}).

\begin{figure}[!ht]
\centering
\includegraphics[width=2.8in]{VE2}
\caption{Valor de $V_E$ en CD --- Parte~2 (NI Multisim).}
\label{fig:ve2}
\end{figure}

\begin{figure}[!ht]
\centering
\includegraphics[width=2.8in]{VC2}
\caption{Valor de $V_C$ en CD --- Parte~2 (NI Multisim).}
\label{fig:vc2}
\end{figure}

\begin{figure}[!ht]
\centering
\includegraphics[width=2.8in]{VCE2}
\caption{Valor de $V_{CE}$ en CD --- Parte~2 (NI Multisim).}
\label{fig:vce2}
\end{figure}

\begin{table}[!ht]
    \centering
    \caption{Par\'ametros CD y CA del amplificador --- Parte~2}
    \begin{tabular}{|>{\centering\arraybackslash}m{1.5cm}|
                     >{\centering\arraybackslash}m{1.5cm}|
                     >{\centering\arraybackslash}m{1.5cm}|
                     >{\centering\arraybackslash}m{1.5cm}|}
        \hline
        Par\'ametro & Calculado & Simulado & Medido \\
        \hline
        $V_B$ (V)            & $1{,}92$ & $1{,}92$ & $1{,}85$ \\
        \hline
        $V_E$ (V)            & $1{,}22$ & $1{,}22$ & $1{,}18$ \\
        \hline
        $I_E$ (mA)           & $2{,}14$ & $2{,}14$ & $2{,}07$ \\
        \hline
        $V_C$ (V)            & $6{,}65$ & $6{,}65$ & $6{,}74$ \\
        \hline
        $V_{CE}$ (V)         & $5{,}43$ & $5{,}43$ & $5{,}58$ \\
        \hline
        $r'_e$ $(\Omega)$    & $11{,}7$ & $11{,}7$ & --- \\
        \hline
        $A_v$                & $-126{,}7$ & $-126{,}7$ & $-112{,}5$ \\
        \hline
        $V_{out}$ (V$_{pp}$) & $2{,}53$ & $2{,}53$ & $2{,}25$ \\
        \hline
    \end{tabular}
    \label{tab:cd_ca_p2}
\end{table}

Los valores de CD deben coincidir con los de la Parte~1, dado que $C_4$, $C_5$ y $C_6$ son circuitos abiertos en CD. La se\~nal de salida simulada se muestra en la Figura~\ref{fig:vout2}.

\begin{figure}[!ht]
\centering
\includegraphics[width=3.2in]{VOUT2}
\caption{Se\~nal de salida $V_{out}$ y entrada $V_s$ simuladas --- Parte~2 (NI Multisim).}
\label{fig:vout2}
\end{figure}


% ---------------------------------------------------------------------------
\subsection{An\'alisis de alta frecuencia}

El an\'alisis utiliza el modelo de Miller simplificado \cite{Floyd2008}. $C_6$ absorbe $C_{be}$, $C_4$ sustituye $C_{bc}$ para el c\'alculo de la capacitancia de Miller, y $C_5$ reemplaza $C_{ce}$.

\begin{equation}
    C_{ent} = C_6 + C_4(|A_v| + 1)
    \label{eq:cent}
\end{equation}

\begin{equation}
    R_{eq(ent)} = (R_A\|R_B)\|R_1\|R_2\|[\beta_{ca}(R_{E1} + r'_e)]
    \label{eq:req_ent}
\end{equation}

\begin{equation}
    f_{c(ent)} = \frac{1}{2\pi \cdot R_{eq(ent)} \cdot C_{ent}}
    \label{eq:fc_ent}
\end{equation}

\begin{equation}
    C_{sal} = C_5 + C_4\frac{|A_v|+1}{|A_v|}
    \label{eq:csal}
\end{equation}

\begin{equation}
    R_c = R_C\|R_L
    \label{eq:rc}
\end{equation}

\begin{equation}
    f_{c(sal)} = \frac{1}{2\pi \cdot R_c \cdot C_{sal}}
    \label{eq:fc_sal}
\end{equation}

\begin{equation}
    f_{cu} = \frac{f_{c(ent)} \cdot f_{c(sal)}}{f_{c(ent)} + f_{c(sal)}}
    \label{eq:fcu}
\end{equation}

\begin{table}[!ht]
    \centering
    \caption{Resultados del an\'alisis de alta frecuencia}
    \begin{tabular}{|>{\centering\arraybackslash}m{0.7cm}|
                     >{\centering\arraybackslash}m{2.1cm}|
                     >{\centering\arraybackslash}m{1.9cm}|
                     >{\centering\arraybackslash}m{1.9cm}|}
        \hline
        Paso & Par\'ametro & Calculado & Medido \\
        \hline
        4  & $C_{ent}$ (nF)     & $12{,}87$ & --- \\
        \hline
        5  & $R_{eq(ent)}$ (k$\Omega$) & $44{,}21$ & --- \\
        \hline
        6  & $f_{c(ent)}$ (kHz)  & $279{,}7$ & --- \\
        \hline
        7  & $C_{sal}$ (pF)     & $200{,}79$ & --- \\
        \hline
        8  & $R_c$ ($\Omega$)         & $2{,}806\,\text{k}\Omega$ & --- \\
        \hline
        9  & $f_{c(sal)}$ (kHz)  & $282{,}51$ & --- \\
        \hline
        10 & $f_{cu}$ (kHz)      & $140{,}55$ & $139{,}8$ \\
        \hline
    \end{tabular}
    \label{tab:alta_freq}
\end{table}


% ---------------------------------------------------------------------------
\subsection{Preguntas de la Parte 2}

\begin{enumerate}

\item \textbf{¿C\'omo afecta la frecuencia cr\'itica superior el resistor sin acople de $10\,\Omega$ en el emisor?}

El resistor $R_{E1} = 10\,\Omega$ sin capacitor de acople incrementa la resistencia de entrada en CA vista desde la base:
\begin{equation}
    Z_{ent(base)} = \beta_{ca}(R_{E1} + r'_e)
\end{equation}
Si $R_{E1}$ estuviera puenteado, este t\'ermino se reducir\'ia a $\beta_{ca} \cdot r'_e$. Sin embargo, al incorporarlo en el paralelo de la ecuaci\'on~(\ref{eq:req_ent}), el resultado est\'a dominado por $R_A\|R_B$, mucho menor que cualquier otra rama. Por tanto, el efecto de $R_{E1}$ sobre $f_{c(ent)}$ es pr\'acticamente despreciable en este circuito particular. En configuraciones con $R_A$ y $R_B$ mayores, su influencia ser\'ia m\'as pronunciada \cite{Floyd2008}.

\item \textbf{¿C\'omo se podr\'ia medir el tiempo de subida para determinar la frecuencia cr\'itica superior?}

Se aplica una se\~nal de onda cuadrada al amplificador en banda media. Con el osciloscopio se mide el tiempo entre el $10\,\%$ y el $90\,\%$ del flanco de subida de la se\~nal de salida. Con ese valor se estima el ancho de banda:
\begin{equation}
    BW = \frac{0{,}35}{t_r}
    \label{eq:bw}
\end{equation}
Dado que $f_{cl} \ll f_{cu}$, se cumple $f_{cu} \approx BW$ \cite{Floyd2008}.

\item \textbf{Usando los resultados experimentales, calcule el tiempo de subida.}

A partir del valor medido de $f_{cu}$:
\begin{equation}
    t_r = \frac{0{,}35}{f_{cu}}
    \label{eq:tr}
\end{equation}

\end{enumerate}


% ===========================================================================
\section{Conclusiones}

El analisis de alta frecuencia evidenció que el efecto Miller incrementa de forma considerable la capacitancia equivalente de entrada del amplificador de manera que reduce el ancho de banda disponible. La comparación entre los parámetros teóricos y experimentales demostró que hubo factores que impactaron en la cercanía de los datos, ya sea por tolerancias de los componentes, capacitancias parásitas del protoboard, y malas conexiones en el circuito físico.

Finalmente, se verificó que el ancho de banda del amplificador depende de la relación entre la ganancia y la respuesta en frecuencia; el aumento en la ganancia de voltaje incrementó el efecto Miller y redujo la frecuencia crítica superior evidenciado ese hecho con su inversa proporcionalidad. Aunque en algunos casos los errores fueron un poco más altos de lo esperado, en general los resultados confirman los conceptos teóricos.

% ===========================================================================
\renewcommand\refname{Referencias}
\begin{thebibliography}{99}
\bibitem{Floyd2008}
T. L. Floyd, \textit{Dispositivos electr\'onicos}, 8\textordfeminine~ed. M\'exico: Pearson Educaci\'on, 2008.

\end{thebibliography}


% ===========================================================================
\appendices
\renewcommand\appendixname{Ap\'endice}
\section{Ecuaciones empleadas --- Parte 1}
\label{apx:ec_p1}

\begin{equation*}
    V_B = V_{CC} \cdot \frac{R_2}{R_1 + R_2}, \quad
    V_E = V_B - V_{BE}, \quad
    I_E = \frac{V_E}{R_{E1} + R_{E2}}
\end{equation*}

\begin{equation*}
    V_C = V_{CC} - I_E R_C, \quad
    V_{CE} = V_C - V_E, \quad
    r'_e = \frac{25\,\text{mV}}{I_E}
\end{equation*}

\begin{equation*}
    A_v = -\frac{R_C\|R_L}{R_{E1} + r'_e}
\end{equation*}

\begin{equation*}
    R_{eq}(C_1) = \beta_{ca}(R_{E1} + r'_e)\|R_1\|R_2 + (R_A + R_{th})\|R_B
\end{equation*}

\begin{equation*}
    R_{umbral} = R_A\|R_B\|R_1\|R_2, \quad
    R_{ent(emisor)} = r'_e + R_{E1} + \frac{R_{umbral}}{\beta_{ca}}
\end{equation*}

\begin{equation*}
    R_{eq}(C_2) = R_{ent(emisor)}\|R_{E2}, \quad
    R_{eq}(C_3) = R_C + R_L
\end{equation*}

\begin{equation*}
    C_{2,\text{nuevo}} = \frac{1}{2\pi \cdot R_{eq}(C_2) \cdot (f_{cl,obj} - f_{C_1} - f_{C_3})}
\end{equation*}


\section{Ecuaciones empleadas --- Parte 2}
\label{apx:ec_p2}

\begin{equation*}
    C_{ent} = C_6 + C_4(|A_v| + 1)
\end{equation*}

\begin{equation*}
    R_{eq(ent)} = (R_A\|R_B)\|R_1\|R_2\|[\beta_{ca}(R_{E1} + r'_e)]
\end{equation*}

\begin{equation*}
    f_{c(ent)} = \frac{1}{2\pi \cdot R_{eq(ent)} \cdot C_{ent}}, \quad
    C_{sal} = C_5 + C_4\frac{|A_v|+1}{|A_v|}
\end{equation*}

\begin{equation*}
    R_c = R_C\|R_L, \quad
    f_{c(sal)} = \frac{1}{2\pi \cdot R_c \cdot C_{sal}}
\end{equation*}

\begin{equation*}
    f_{cu} = \frac{f_{c(ent)} \cdot f_{c(sal)}}{f_{c(ent)} + f_{c(sal)}}, \quad
    t_r = \frac{0{,}35}{f_{cu}}
\end{equation*}


\section{Informaci\'on del autor}

\begin{IEEEbiographynophoto}{Brandon N\'ajera Zu\~niga}
Estudiante de Ingenier\'ia Electr\'onica del Instituto Tecnol\'ogico de Costa Rica. Cursando el curso EL3215 Laboratorio de Electr\'onica Anal\'ogica, I Semestre 2026.

Correo: brandonnz@estudiantec.cr
\end{IEEEbiographynophoto}
\begin{IEEEbiographynophoto}{Kevin Mora Sobalvarro}
Estudiante de Ingenier\'ia Electr\'onica del Instituto Tecnol\'ogico de Costa Rica. Cursando el curso EL3215 Laboratorio de Electr\'onica Anal\'ogica, I Semestre 2026.

Correo: kevinmorasoval03@estudiantec.cr
\end{IEEEbiographynophoto}
\vfill

\end{document}